Khóa luận này đã xây dựng một nội dung hoàn chỉnh từ nền tảng toán học
đến phân tích lý thuyết và kiểm chứng thực nghiệm của thuật toán
Deutsch--Jozsa. Cụ thể, Chương~1 thiết lập cơ sở toán học cần thiết gồm
không gian Hilbert, các toán tử Hermitian và unitary, Định lý phổ, tích
tensor và bốn tiên đề Dirac--von Neumann, trong đó cơ chế giao thoa lượng
tử được phân tích từ nền tảng tiên đề và theo dõi xuyên suốt đến bước cuối
của thuật toán. Chương~2 triển khai cơ sở toán học thành mô hình tính toán
lượng tử cụ thể qua qubit, các cổng lượng tử cơ bản và cấu trúc mạch.
Chương~3 phát biểu bài toán Deutsch--Jozsa, chứng minh chặt chẽ kết quả
trung tâm rằng thuật toán phân biệt hàm hằng và hàm cân bằng với xác suất
$1$ trong đúng một lần truy vấn oracle, và kiểm chứng kết quả đó qua mô
phỏng số với Qiskit cho $n = 1, \ldots, 5$.

Bên cạnh những kết quả đạt được, khóa luận còn một số hạn chế cần thừa
nhận. Về phạm vi thuật toán, Deutsch--Jozsa giải quyết một bài toán có điều
kiện và ưu thế tuyệt đối $O(1)$ so với $O(2^{n-1})$ chỉ đúng với thuật
toán cổ điển tất định; nếu cho phép xác suất, ưu thế này thu hẹp đáng kể.

Trong tương lai, đề tài có thể được phát triển theo hướng mở rộng sang các
thuật toán lượng tử có ứng dụng thực tiễn hơn. Cụ thể, thuật toán Grover
khai thác cùng cơ chế giao thoa để tìm kiếm với độ phức tạp $O(\sqrt{N})$,
biến đổi Fourier lượng tử (QFT) tổng quát hóa Hadamard transform đã xây
dựng trong khóa luận, và thuật toán Shor sử dụng QFT để phân tích nhân tử
với hiệu quả vượt trội so với mọi thuật toán cổ điển đã biết.